\section*{Chapitre \ref{ch:iv2:making-a-market} --- Market making}
\markboth{\protect\omTocMark{}{Réponses~: Market making}}{\protect\omTocMark{}{Réponses~: Market making}}

\begin{solution}{iq:iv2:making-a-market:1}
Coter 10 au cours acheteur (bid), 11 au cours vendeur (ask), et espérer gagner 0.5. Une cotation de 11 au bid, 12 à l'ask rapporte également 0.5 en moyenne.

Chaque dé vaut 3.5 en moyenne, donc la somme vaut 10.5. Si le client vous vend à 10, vous gagnez 0.5 en moyenne, et s'il vous achète à 11, vous gagnez 0.5. À 11 et 12, acheter à 11 fait perdre 0.5 et vendre à 12 rapporte 1.5, et la moyenne des deux vaut à nouveau 0.5. Un client dont le sens est à pile ou face vous paie la moitié de votre largeur de la fourchette, quelle que soit la position de la cotation.

\iqwrong{Conclure que la position de la cotation n'a pas d'importance. Elle n'en a pas seulement parce que ce client l'ignore. Quelqu'un qui a calculé 10.5 vous vendrait à 11 à chaque fois, et vous perdriez 0.5 à chaque transaction.}
\iqfollow{Quel risque prenez-vous pour ces 0.5~?}{La somme a un écart type de $\sqrt{35/4} \approx 2.96$, soit environ six fois votre avantage. Une seule transaction ne prouve rien. L'avantage n'apparaît que sur un grand nombre de transactions.}
\iqwhy{La largeur vous rémunère pour les flux non informés, et le centre vous protège contre tous les autres. Dites les deux en donnant une cotation.}
\end{solution}

\begin{solution}{iq:iv2:making-a-market:2}
Coter 6 au bid, 8 à l'ask, et perdre $30/13 \approx 2.31$ en moyenne. La plus petite largeur entière qui perd moins de 1 est 6, avec 4 au bid et 10 à l'ask, qui perd $12/13 \approx 0.92$.

Les rangs sont équiprobables et centrés sur 7, donc la cotation symétrique est la meilleure. À 6 et 8, les rangs 1 à 5 vous vendent et vous font perdre 5, 4, 3, 2 et 1, et les rangs 9 à 13 vous achètent et vous font perdre 1 à 5. Le total fait 30 sur 13 rangs. La meilleure cotation à chaque largeur perd~:
\begin{center}\small
\begin{tabular}{@{}lccccc@{}}
Largeur & 2 & 4 & 5 & 6 & 12 \\ \hline
Perte moyenne & 2.31 & 1.54 & 1.23 & 0.92 & 0
\end{tabular}
\end{center}

\iqwrong{Chercher une cotation qui rapporte de l'argent. Quelqu'un qui connaît le rang ne trade que lorsque vous avez tort, donc vous ne pouvez jamais profiter. Ce n'est qu'à 1 au bid et 13 à l'ask que vous cessez de perdre, et là personne ne trade. La largeur est un prix que vous fixez pour vous protéger d'être pris à contre-pied.}
\iqwhy{C'est le cas extrême de la sélection adverse\index{sélection adverse}. Quand vous soupçonnez que l'autre côté est informé, la question cesse d'être «~combien cela vaut-il~?~» pour devenir «~combien suis-je prêt à payer pour le savoir~?~»}
\end{solution}

\begin{solution}{iq:iv2:making-a-market:3}
Coter 14 au bid, 16 à l'ask. Le contrat vaut $91/6 \approx 15.17$, et non $3.5^2 = 12.25$.

L'espérance vaut $(1 + 4 + 9 + 16 + 25 + 36)/6 = 91/6$. Parmi les demi-pas, un milieu de 15 est le plus proche de 15.17 qui tienne dans une largeur de 2, donc la cotation va de 14 à 16. Le sens du client est à pile ou face, donc l'avantage vaut 1 en moyenne.

\iqwrong{Élever au carré le dé moyen. L'espérance d'un carré est le carré de l'espérance plus la variance, et la variance du dé vaut $35/12 \approx 2.92$, ce qui correspond exactement à l'écart entre 12.25 et 15.17. C'est l'inégalité de Jensen\index{inégalité de Jensen}~: un gain convexe vaut plus que le gain évalué à la moyenne.}
\iqfollow{Et un contrat qui paie 1 divisé par le dé~?}{Il est également convexe en le dé, donc là encore le gain moyen dépasse le gain évalué à la moyenne~: $49/120 \approx 0.41$ contre $1/3.5 \approx 0.29$.}
\end{solution}

\begin{solution}{iq:iv2:making-a-market:4}
Coter la prochaine carte 5.75 au bid, 6.75 à l'ask, et la somme des autres 263.5 au bid, 264.5 à l'ask. La première porte un risque réel, et la seconde n'en porte aucun.

Chaque rang de 1 à 13 apparaît quatre fois, donc le jeu entier totalise $4 \times 91 = 364$. Avec 100 visibles, les 42 cartes non vues totalisent exactement 264. Chacune a la même probabilité d'être la suivante, donc la prochaine carte vaut $264/42 = 44/7 \approx 6.29$, et le quart le plus proche donne un milieu de 6.25. La valeur de votre contrepartie se situe à l'intérieur des deux cotations, donc elle accepte.

\iqwrong{Traiter la somme comme incertaine parce qu'elle additionne 42 cartes aléatoires. Ces cartes ne sont pas indépendantes. Le total est fixé, donc la variance de la somme est nulle. Calculer la moyenne $42 \times 7 = 294$ ignore ce que vous avez vu.}
\iqfollow{Que feriez-vous si quelqu'un vous proposait de vendre la somme à 263~?}{Acheter autant qu'il en vendra, car vous gagnez 1 sur chaque lot avec certitude. Quiconque sait compter le jeu ne le proposera pas.}
\end{solution}

\begin{solution}{iq:iv2:making-a-market:5}
Ils vous vendent à la fois $A$ à 0.55 et $B$ à 0.50, en encaissant 1.05 et en vous devant exactement 1, soit un gain certain de 0.05. Pour l'exclure, les deux cours acheteurs doivent totaliser au plus 1, et les deux cours vendeurs au moins 1.

Exactement un de $A$ et $B$ paie, car les deux événements couvrent tous les résultats du dé sans recouvrement, donc une paire d'un lot chacun vaut exactement 1 quoi qu'il arrive. Une paire qu'on vous vend plus de 1 est un gain gratuit pour l'autre côté. Une paire que vous achetez à moins de 1 est aussi un gain gratuit. Vos cours vendeurs totalisent ici 1.12, donc ce côté est sûr.

\iqfollow{Est-ce sûr de coter 0.40 à 0.45 sur les deux~?}{Non. Les cours vendeurs totalisent 0.90, donc ils vous achètent les deux pour 0.90 et reçoivent 1, soit 0.10 de gain certain.}
\end{solution}

\begin{solution}{iq:iv2:making-a-market:6}
Cotez 4 au bid, 5 à l'ask avant les lancers, et 3.5 au bid, 4.5 à l'ask après le 3. Après un 6, le contrat vaut exactement 6, donc vous afficheriez 5.5 au bid, 6.5 à l'ask.

Le plus grand de deux dés vaut $k$ dans $2k - 1$ des 36 issues, donc la valeur vaut $\sum k(2k - 1)/36 = 161/36 \approx 4.47$. Sachant que le premier dé vaut 3, le second dé donne un maximum de 3 dans trois cas, et sinon 4, 5 ou 6, donc la valeur vaut $(3 + 3 + 3 + 4 + 5 + 6)/6 = 4$.

Après le premier dé, les valeurs sont~:
\begin{center}\small
\begin{tabular}{@{}l*{6}{c}@{}}
Premier dé & 1 & 2 & 3 & 4 & 5 & 6 \\ \hline
Valeur & 3.50 & 3.67 & 4.00 & 4.50 & 5.17 & 6.00
\end{tabular}
\end{center}

\iqfollow{Pourquoi garderiez-vous jamais une largeur après le 6~?}{Vous n'en avez pas besoin pour le risque, car rien n'est incertain. L'énoncé vous oblige à en afficher une, et le pile ou face du client vous la fait gagner. Un vrai client ne traderait pas à 6.5 pour quelque chose qui vaut 6.}
\end{solution}

\begin{solution}{iq:iv2:making-a-market:7}
Coter environ 159 au bid, 179 à l'ask. Il y a 168 nombres premiers inférieurs à 1\,000, donc cette cotation passe.

Le théorème des nombres premiers\index{théorème des nombres premiers} dit que le nombre de premiers jusqu'à $n$ vaut environ $n/\ln n$, et une forme plus fine vaut environ $n/(\ln n - 1)$. Avec $\ln 1000 = 3 \times 2.303 \approx 6.91$, la forme plus fine donne $1000/5.91 \approx 169$. Le centre de la cotation vaut 169, et une largeur de 20 la fait aller de 159 à 179.

Si vous ne vous rappelez pas le théorème, criblez. Environ $1000 \times \tfrac12 \times \tfrac23 \times \tfrac45 \times \tfrac67 \approx 229$ nombres n'ont aucun facteur 2, 3, 5 ou 7, et le compte exact est 228. Ils comprennent encore des composés comme $11 \times 11$, donc la réponse est bien en dessous de 229.

\iqwrong{Utiliser $n/\ln n \approx 145$ tel quel. La cotation de 135 à 155 perd alors $168 - 155 = 13$ lorsque l'interrogateur achète à 155. La forme simple sous-estime pour les petits $n$.}
\end{solution}

\begin{solution}{iq:iv2:making-a-market:8}
Cotez 5 au bid, 7 à l'ask. Les deux côtés gagnent 1 en moyenne, et ils diffèrent entièrement par leur forme. Si vous êtes long à 5, vous perdez de l'argent avec une probabilité de $1 - (5/6)^4 \approx 52~\%$, et au plus 4. Si vous êtes short à 7, vous perdez avec une probabilité de $(5/6)^7 \approx 28~\%$, et la perte n'a pas de limite.

Le nombre de lancers suit une loi géométrique\index{loi géométrique} de moyenne 6 et de variance 30, donc l'écart type vaut environ 5.5. Sa médiane vaut 4, car $1 - (5/6)^3$ est inférieur à un demi et $1 - (5/6)^4$ est supérieur. La plupart des résultats sont inférieurs à la moyenne, et quelques-uns sont très au-dessus.

Une position longue perd donc en général un peu et gagne parfois beaucoup, et une position courte gagne en général un peu et perd parfois beaucoup. Il faut 20 lancers ou plus environ 3~\% du temps, et alors un short perd au moins 13.

\iqwrong{Centrer sur la médiane parce que «~la plupart des issues sont inférieures à 6~». Le contrat paie le nombre de lancers, donc c'est la moyenne qui permet d'équilibrer les gains sur un grand nombre de transactions.}
\iqfollow{Quel côté préféreriez-vous avec le même gain moyen~?}{Cela dépend de vos limites. Un short a une petite probabilité d'une très grosse perte, il faut donc une limite, plus d'avantage, ou les deux. Dites-le avant de dire lequel vous préférez.}
\iqwhy{Vendre un gain asymétrique à une moyenne équitable, c'est ainsi que l'on vend des options. Une position peut avoir une bonne espérance et rester une position que vous ne pouvez pas vous permettre de tenir.}
\end{solution}

\begin{solution}{iq:iv2:making-a-market:9}
Cotez 4 au bid, 6 à l'ask. Le résultat tombe hors de la cotation avec une probabilité d'environ 29~\%. Après une première carte rouge, la valeur vaut $92/17 \approx 5.41$, et vous cotez 4.5 au bid, 6.5 à l'ask.

Par symétrie, chaque carte est rouge avec probabilité un demi, donc le nombre moyen de rouges vaut 5 quelle que soit la dépendance entre les cartes. Le compte suit une loi hypergéométrique\index{loi hypergéométrique}, de variance
\[
10 \cdot \tfrac12 \cdot \tfrac12 \cdot \tfrac{52 - 10}{52 - 1} = \tfrac{35}{17},
\]
donc l'écart type vaut 1.43, en dessous des 1.58 de dix lancers de pièce indépendants. Tirer sans remise ramène le compte vers le milieu. En sommant les probabilités exactes de 4, 5 et 6, on obtient 0.71.

Après un rouge, les neuf autres cartes proviennent de 25 rouges sur 51, donc la valeur vaut $1 + 9 \cdot 25/51$. La première carte valait 0.5 et vaut maintenant 1, soit un gain de 0.5. Les neuf autres passent de 4.5 à 4.41, car une rouge a quitté le jeu. Au total, la valeur monte de 0.41.

\iqwrong{Monter de 1, ou de 0.5. Une hausse de 1 oublie que la première carte comptait déjà pour 0.5. Une hausse de 0.5 oublie que les neuf autres cartes dépendent de la première.}
\end{solution}

\begin{solution}{iq:iv2:making-a-market:10}
Avec $k = h/\sigma$, la perte espérée vaut $2\sigma\left[\varphi(k) - k\,\Phi(-k)\right]$, où $\varphi$ et $\Phi$ sont la densité et la fonction de répartition de la loi normale centrée réduite. Un intervalle à 90~\%, $k = 1.645$, perd $0.042\,\sigma$. Pour perdre moins de $0.05\,\sigma$, il faut $k \ge 1.57$, soit une largeur d'environ $3.1\,\sigma$.

La contrepartie achète lorsque $X > \mu + h$, ce qui vous coûte $X - \mu - h$, et vend dans le cas symétrique. La perte vaut donc
\[
2\,\mathbb{E}\left[(X - \mu - h)^+\right] = 2\sigma\int_k^\infty (z - k)\,\varphi(z)\,dz = 2\sigma\left[\varphi(k) - k\,\Phi(-k)\right].
\]
C'est la fonction de perte normale\index{fonction de perte normale}. En $k = 0$ elle vaut $0.80\,\sigma$, et en $k = 1$ elle vaut $0.167\,\sigma$.

\iqwrong{Penser qu'un intervalle à 90~\% perd 10~\% de quelque chose. Vous négociez 10~\% du temps, mais ces transactions ne perdent qu'une petite somme en moyenne, car la queue est mince juste au-delà de l'intervalle.}
\iqfollow{Pourquoi ne pas coter très large, alors~?}{Face à cette contrepartie, on doit le faire. Mais on veut aussi le flux non informé, et celui-ci cesse de trader si l'on est large. Jusqu'où l'on va est un compromis entre les deux, et non la simple expression de son incertitude.}
\end{solution}

\begin{solution}{iq:iv2:making-a-market:11}
Centrer sur environ 254. Le contrat verse le nombre, donc la fair value est la moyenne, et pour une loi lognormale\index{loi lognormale} la moyenne vaut la médiane multipliée par $e^{\sigma^2/2}$. Ici $\sigma = \ln 2$, donc le facteur vaut $e^{0.24} \approx 1.27$.

Donc $200 \times 1.27 \approx 254$. Le vrai nombre dépasse 254 avec une probabilité $\Phi(-\sigma/2) = \Phi(-0.35) \approx 36~\%$. L'intervalle à 90~\% est $200 \cdot 2^{\pm 1.645}$, soit environ 64 à 625.

\iqwrong{Centrer sur la meilleure estimation de 200. La moitié du temps, le nombre est au-dessus de 200, mais il peut être bien plus haut, alors qu'il ne peut être qu'un peu plus bas en valeur absolue. Si l'on vend à 200 encore et encore, on perd beaucoup, 54 en moyenne. Une cotation centrée sur la médiane se trompe aussi sur la forme~: l'intervalle n'est pas symétrique autour du centre.}
\iqfollow{Comment avez-vous obtenu un facteur de 2~?}{Une estimation de Fermi est un produit de facteurs, et les logarithmes d'un produit s'additionnent. Trois facteurs, chacun incertain d'un facteur 1.5, donnent un écart-type logarithmique de $\sqrt3 \ln 1.5$, ce qui fait un facteur d'environ 2.0 au total.}
\iqwhy{Les quantités inconnues sont en général des produits, donc leur incertitude est multiplicative. Votre cours acheteur peut être bien plus proche de zéro que votre cours vendeur ne l'est de votre centre, et le dire est la marque de quelqu'un qui y a réfléchi.}
\end{solution}

\begin{solution}{iq:iv2:making-a-market:12}
Coter un cours acheteur de 94 et un cours vendeur de 106. Le cours vendeur d'équilibre est l'espérance de la valeur sachant qu'on achète, et le cours acheteur d'équilibre est l'espérance de la valeur sachant qu'on vend.

Appliquer la formule de Bayes\index{formule de Bayes}. Si la valeur vaut 120, un achat survient avec une probabilité $0.30 + 0.70 \cdot \tfrac12 = 0.65$. Si elle vaut 80, un achat ne peut venir que du tirage à pile ou face, $0.70 \cdot \tfrac12 = 0.35$. Donc, sachant qu'il y a un achat, la probabilité que la valeur soit 120 vaut $0.65/(0.65 + 0.35) = 0.65$, et le cours vendeur vaut $80 + 40 \times 0.65 = 106$. Le cours acheteur en est l'image symétrique, 94.

Donc la largeur de la fourchette est la part des traders informés multipliée par l'écart entre les valeurs, $0.30 \times 40 = 12$. C'est la cotation du modèle de Glosten--Milgrom\index{modèle de Glosten--Milgrom}. Les traders non informés paient exactement ce que les informés vous prennent.

\iqwrong{Coter autour de 100 avec une largeur de fourchette choisie au jugé. Ce qui compte est la valeur sachant qu'on achète. Il existe ici un raccourci~: 30~\% des acheteurs sont informés et paient 120, et les autres paient en moyenne 100, donc le cours vendeur vaut $0.3 \times 120 + 0.7 \times 100 = 106$. Cela ne marche que parce que la probabilité a priori vaut un demi. Avec une autre probabilité a priori, la part des informés parmi les acheteurs diffère de 30~\%.}
\iqfollow{Et si tout le monde était informé, ou si personne ne l'était~?}{Si tout le monde est informé, la cotation est de 80 et 120, et le trader ne gagne rien à ces prix. Si personne ne l'est, elle est de 100 et 100, sans fourchette du tout, puisqu'il n'y a rien contre quoi se protéger dans ce modèle.}
\end{solution}

\begin{solution}{iq:iv2:making-a-market:13}
Les cotations précédant chacun des cinq traders figurent ci-dessous, et l'ordre n'a pas d'importance, car chaque achat multiplie la cote de 100 par 1.5, et chaque vente la divise par 1.5.
\begin{center}\small
\begin{tabular}{@{}lccccc@{}}
Trader & 1 & 2 & 3 & 4 & 5 \\ \hline
Transactions jusqu'ici & aucune & B & BB & BBS & BBSB \\
Cours acheteur & 40.00 & 50.00 & 60.00 & 50.00 & 60.00 \\
Cours vendeur & 60.00 & 69.23 & 77.14 & 69.23 & 77.14
\end{tabular}
\end{center}

Si la valeur vaut 100, un achat a une probabilité $0.20 + 0.40 = 0.60$, et si elle vaut 0, il a une probabilité $0.40$. Le rapport de vraisemblance\index{rapport de vraisemblance} vaut 1.5, et une vente donne son inverse. Ainsi, avec $n$ achats nets, la cote de 100 vaut $1.5^n$. Le cours vendeur correspond à la valeur sachant un achat de plus, à la cote $1.5^{n+1}$, et le cours acheteur à la valeur sachant une vente de plus, à la cote $1.5^{n-1}$. Avec $n = 1$, le cours vendeur vaut $2.25/3.25$ de 100, soit 69.23. La cote finale vaut $1.5^2$, donc la probabilité de 100 vaut $9/13$.

\iqfollow{Et si les traders non informés achetaient 70~\% du temps~?}{Un achat a maintenant une probabilité de $0.76$ ou $0.56$, soit un rapport de 1.36, et une vente a un rapport $0.24/0.44 = 0.55$. Comme $1/1.36 = 0.74$, une vente fait davantage bouger la cote qu'un achat~: la transaction la plus rare en dit davantage. L'ordre n'a toujours pas d'importance, mais le solde net ne suffit plus, et il faut les deux comptages.}
\end{solution}

\begin{solution}{iq:iv2:making-a-market:14}
Le contrat vaut 2.7 après la vente, et vous cotez 2.2 au cours acheteur et 3.2 au cours vendeur. Après une passe, il vaut 3.5.

Pondérer chaque valeur par la probabilité d'une vente. Pour les valeurs 1 et 2, l'observateur vend, donc la probabilité vaut $\tfrac12 + \tfrac12 \cdot \tfrac12 = \tfrac34$. Pour les valeurs 3 à 6, seul le tirage à pile ou face peut vendre, donc la probabilité vaut $\tfrac14$. Les poids somment à $\tfrac52$, donc la loi a posteriori vaut $0.3, 0.3, 0.1, 0.1, 0.1, 0.1$, et la moyenne vaut 2.7.

Une passe ne peut venir que d'un observateur, et un observateur ne passe que sur 3 ou 4. La valeur est donc l'une de ces deux, de moyenne 3.5. Notons que cela laisse le mid là où il était, tandis que l'écart s'est réduit à une seule unité.

Pour vérifier, appliquer la loi de l'espérance totale\index{loi de l'espérance totale}. Une vente et un achat ont chacun une probabilité $\tfrac{5}{12}$, une passe a une probabilité $\tfrac16$, et la valeur sachant un achat vaut 4.3 par symétrie. Alors $\tfrac{5}{12} \cdot 2.7 + \tfrac{5}{12} \cdot 4.3 + \tfrac16 \cdot 3.5 = 3.5$, soit la valeur initiale.

\iqwrong{Supposer que le vendeur a forcément observé, ce qui donne 1.5. Le vendeur peut très bien être non informé, et c'est ce qui empêche la mise à jour d'aller si loin.}
\end{solution}

\begin{solution}{iq:iv2:making-a-market:15}
Le cours vendeur est d'environ 100.23 pour un lot et 106.40 pour 100 lots. Un cours vendeur unique pour les deux tailles est d'environ 106.0.

La valeur vaut 110 si l'acheteur est informé, et elle vaut en moyenne 100 s'il ne l'est pas. Le cours vendeur vaut donc $100 + 10 \times$ la probabilité que l'acheteur soit informé, sachant la taille. La taille est le signal, et la formule de Bayes donne
\[
\Pr(\text{informé} \mid 100) = \frac{0.1 \cdot 0.8}{0.1 \cdot 0.8 + 0.9 \cdot 0.05} = 0.64, \qquad
\Pr(\text{informé} \mid 1) = \frac{0.1 \cdot 0.2}{0.1 \cdot 0.2 + 0.9 \cdot 0.95} \approx 0.023.
\]

Un cours vendeur unique doit être à l'équilibre sur l'ensemble des lots échangés, et non sur les traders. Les traders informés envoient $0.1(0.2 + 80) = 8.02$ lots par arrivée, et les autres en envoient $0.9(0.95 + 5) = 5.355$, donc les traders informés envoient 60~\% des lots, et le cours vendeur vaut environ 106.0.

\iqwrong{Faire la moyenne sur les traders. La plupart des traders envoient un lot et ne sont pas informés, mais la plupart des lots, eux, le sont. Ce sont ces lots qui vous coûtent de l'argent.}
\iqfollow{Que fait le cours vendeur unique aux petits clients~?}{Ils paient 106 pour quelque chose qui vaut à peine plus de 100, donc ils partent. Et un trader informé, dont le prix dépend de la taille, se fractionnera en petits ordres, ce qui explique que le signal s'affaiblisse avec le temps.}
\iqwhy{Afficher moins de taille près du mid et davantage plus loin découle de ce calcul.}
\end{solution}

\begin{solution}{iq:iv2:making-a-market:16}
Le cours vendeur le plus bas est 99.21 et le cours acheteur le plus haut est 99.19, donc le centre vaut 99.20, et vous affichez 98.20 au cours acheteur et 100.20 au cours vendeur. En position courte de 40, vous affichez 99.80 au cours acheteur et 101.80 au cours vendeur. Le centre se déplace de 0.02 à l'encontre de chaque transaction.

Vendre un lot à $a$ depuis une position longue de 40 modifie votre score de $a - 100 + 0.01(40^2 - 39^2) = a - 100 + 0.79$, donc vous acceptez tout prix $a \ge 99.21$. Acheter un lot à $b$ modifie votre score de $100 - b - 0.01(41^2 - 40^2) = 100 - b - 0.81$, donc vous acceptez tout prix $b \le 99.19$.

Avec une position $q$, le centre vaut $100 - 2 \times 0.01 \times q$. Les prix pour lesquels vous êtes indifférent sont ce centre plus ou moins 0.01. Un achat fait monter $q$ de 1, et le centre baisser de 0.02. Ce déplacement simultané des deux côtés est un skew d'inventaire\index{skew d'inventaire}.

\iqwrong{Baisser seulement le cours acheteur. Vous voulez réduire la position dans les deux sens~: être moins susceptible d'acheter davantage, et plus susceptible de vendre. Déplacer les deux côtés conserve la largeur de votre fourchette et affiche le même skew.}
\iqfollow{Pourquoi un score comporterait-il une position au carré~?}{C'est le coût de portage du risque. Une position de $q$ sur un prix de variance $\sigma^2$ a une variance de $q^2\sigma^2$, et un détenteur averse au risque facture proportionnellement à celle-ci.}
\iqwhy{Beaucoup de recruteurs posent cette question juste après une transaction~: «~vous êtes maintenant long, recotez-moi.~» Un candidat qui déplace la fair value confond ce qu'il détient avec ce que la chose vaut.}
\end{solution}

\begin{solution}{iq:iv2:making-a-market:17}
Vous gagnez en espérance 9.125 sur 10 manches, contre 10 sans limite. À long terme, vous gagnez $6/7 \approx 0.86$ par manche. Avec une limite de $\pm L$, cela vaut $2L/(2L + 1)$ par manche.

Chaque transaction rapporte 1, donc la question est de savoir combien de clients vous refusez. Une position à la limite en refuse la moitié. En partant d'une position neutre, vous ne pouvez refuser un client pour la première fois qu'à la manche 4, et en faisant avancer exactement les probabilités de la position, on obtient $73/8$ transactions sur 10 manches.

À long terme, la position est une marche aléatoire qui reste sur place lorsqu'elle franchirait une limite. Chaque ligne et chaque colonne de sa matrice de transition somment à 1, donc sa loi stationnaire\index{loi stationnaire} est uniforme sur les $2L + 1$ positions. Vous êtes donc à une limite avec une probabilité $2/(2L + 1)$, et vous refusez alors la moitié des clients, soit une part de $1/(2L + 1)$. Pour $L = 3$, cela fait un client sur sept.

\iqwrong{Supposer que la limite ne pèse pas, en se disant «~nous ne sommes à la limite que 2 fois sur 7~». Cela coûte un septième de votre revenu, et ce, avec un flux entièrement bilatéral.}
\iqfollow{Comment feriez-vous mieux~?}{On pratique un skew avant d'atteindre une limite, de sorte que le flux ramène à une position neutre. Ici, les clients ignorent le prix, donc le skew ne fait rien. De vrais clients, eux, ne l'ignorent pas.}
\end{solution}

\begin{solution}{iq:iv2:making-a-market:18}
Les contreparties peuvent verrouiller un profit si votre cours acheteur est supérieur à 10.8 ou si votre cours vendeur est inférieur à 10.2. Si vous couvrez la vente à 10.7, vous verrouillez une perte de 0.1. Si vous ne couvrez pas, vous gagnez 0.2 en moyenne, avec un écart-type d'environ 2.96.

Trois contrats à un seul dé versent exactement la même somme que $S$. Les trois contrats coûtent donc 10.8 à l'achat et rapportent 10.2 à la vente. Si vous cotez un cours acheteur supérieur à 10.8, par exemple 10.9, la contrepartie achète les trois morceaux à 10.8 et vous les vend à 10.9, gagnant 0.1 à coup sûr. Si vous cotez un cours vendeur inférieur à 10.2, la contrepartie fait l'inverse. Entre 10.2 et 10.8, vous êtes à l'abri de cet arbitrage.

Un client lève votre cours vendeur de 10.7, et la fair value vaut 10.5, donc vous disposez d'un avantage (edge) de 0.2. Acheter les morceaux à 10.8 fait perdre cet avantage et coûte 0.1 de plus. Une cotation située à l'intérieur de la fourchette propre aux morceaux ne peut donc gagner de l'argent qu'en portant le risque un certain temps~: vous gardez la position jusqu'à l'arrivée d'un client de l'autre côté, ou vous estimez que le risque en vaut 0.2.

\iqwrong{Supposer qu'une couverture est toujours la bonne solution. Une couverture a un prix. Il faut comparer son coût au risque qu'elle supprime.}
\iqfollow{Où coteriez-vous si vous deviez couvrir chaque transaction sur-le-champ~?}{Aux niveaux 10.2 et 10.8, ou en dehors de ces niveaux. Une transaction verrouille alors un gain nul ou meilleur.}
\end{solution}

\begin{solution}{iq:iv2:making-a-market:19}
$A$ vaut désormais $28/3 \approx 9.33$ et $B$ vaut $49/6 \approx 8.17$. La meilleure couverture ne représente que $4/23 \approx 0.17$ d'un lot de $B$, et elle supprime environ 8.7~\% de la variance.

Les contreparties achètent lorsque $X + Y$ est au moins égal à 8, ce qui arrive avec une probabilité $15/36$. Alors $\mathbb{E}[X + Y] = 140/15 = 9.33$. Par symétrie, $X$ prend la moitié de l'excédent au-delà de 7, donc $B$ monte de moitié autant que $A$~: elle passe de 7 à 8.17. La mise à jour croisée\index{ratio de couverture} est la pente de $B$ sur $A$, soit $\tfrac12$.

Pour la couverture, calculer en tenant compte de ce que l'on sait maintenant. Vous êtes en position short sur $A$, et sachant $X + Y \ge 8$, la variance de $A$ vaut $14/9 \approx 1.56$, la variance de $B$ vaut $161/36 \approx 4.47$, et leur covariance vaut $7/9$. La quantité de $B$ qui minimise la variance vaut $(7/9)/(161/36) = 4/23$. Elle supprime $(7/9)^2/\left(\tfrac{14}{9} \cdot \tfrac{161}{36}\right) \approx 8.7~\%$.

\iqwrong{Utiliser le ratio inconditionnel de $\tfrac12$. Avant toute transaction juste, et il supprime 25~\% de la variance de $A$. Or $A$ est désormais cantonnée dans un intervalle étroit, et $B$ porte encore tout le bruit de $Z$. Une couverture de $\tfrac12$ fait passer la variance de 1.56 à 1.90.}
\iqfollow{Et si elles n'avaient connu que $Y$~?}{Alors $B$ ne bouge pas du tout. Le degré de déplacement d'une cotation lorsqu'une autre s'exécute dépend de ce que le trader sait, et il faut le deviner.}
\end{solution}

\begin{solution}{iq:iv2:making-a-market:20}
$h = 1.6/10.8 \approx 0.148$. Avec une probabilité de 50~\% d'être touché, $h \approx 0.095$. Avec 5 clients, $h \approx 0.276$.

Sur un cycle, vous gagnez $10h$ auprès des clients. Si le trader rapide vous touche, vous vendez à $f + h$ un actif qui vaut désormais $f + 2$, ou vous achetez en miroir, et dans les deux cas vous perdez $2 - h$. Le point mort exige donc
\[
n h - p (M - h) = 0, \qquad h = \frac{p M}{n + p},
\]
avec $n = 10$ clients, $p = 0.8$ et un mouvement $M = 2$. Chaque client finit par payer sa part de la nouvelle information, soit environ $pM/n$ lorsque $n$ est grand.

\iqwrong{Comparer $h$ seulement au mouvement de 2, et en conclure qu'il faut une fourchette de 2. Vous n'êtes pas touché sur chaque client, mais une seule fois par information. La largeur de la fourchette est ce coup unique réparti sur l'ensemble des clients.}
\iqfollow{Que dit cela de la vitesse~?}{Elle a une valeur monétaire, et vous pouvez la chiffrer. Ramener $p$ de 0.8 à 0.5 abaisse la demi-fourchette de point mort de 0.053. Vous pouvez donc coter plus serré, prendre davantage de clients, ou garder la différence. C'est aussi pourquoi un teneur de marché plus lent élargit sa cotation, ou cote moins, avant une annonce programmée.}
\iqwhy{C'est le modèle le plus simple de la sélection adverse subie lorsque votre cotation est périmée, un coût que tout desk de cotation doit couvrir.}
\end{solution}

\begin{solution}{iq:iv2:making-a-market:21}
Rejoindre rapporte 0.40 par trader. Se placer à l'intérieur avec une demi-fourchette $h$ rapporte $h - 0.2$, donc $h = 0.5$ rapporte 0.30, et $h = 0.9$ rapporte 0.70. Si le rival reste en place, 0.9 est le meilleur choix. S'il bouge aussi, les cotations descendent à une demi-fourchette de 0.3 à 0.4, juste au-dessus de la demi-fourchette de point mort de 0.2.

Par transaction, à demi-fourchette $h$, un trader non informé vous paie $h$, et un trader informé vous coûte $2 - h$. Ensemble, $0.9h - 0.1(2 - h) = h - 0.2$. À 99 et 101, cela fait 0.8 par transaction, et le partage donne 0.40. Se placer à l'intérieur l'emporte sur rejoindre lorsque $h - 0.2 > 0.4$, soit $h > 0.6$, si bien que les petites améliorations gagnent.

Lorsque les deux bougent, un seul pas capte tout le flux au lieu de la moitié. À demi-fourchette commune $h$, vous gagnez $(h - 0.2)/2$, et passer devant de 0.1 rapporte $h - 0.3$. Passer devant paie lorsque $h > 0.4$. Il y a égalité à 0.4, et à 0.3 c'est perdant. La course s'arrête donc là, à un ou deux pas au-dessus du point mort, à cause de la taille du pas. Au point mort, la largeur vaut $2 \times 0.2 = 0.4$, soit la part informée multipliée par l'intervalle, $0.1 \times 4$.

\iqfollow{Et si une règle permet au premier de la file de prendre toute la transaction~?}{Alors rejoindre vaut moins, voire rien si vous rejoignez derrière. Être en tête de file a une valeur propre, comme sur une bourse réelle avec la priorité prix-temps.}
\end{solution}

\begin{solution}{iq:iv2:making-a-market:22}
$h = 0.8$, soit une largeur de 1.6. Avec un quart, $h = 0.375$. Avec trois quarts, $h = 4/3 \approx 1.33$.

Quelqu'un qui a vu un dé $d$ valorise la somme à $d + 7$, soit de 8 à 13, ou à $\pm 0.5, \pm 1.5, \pm 2.5$ autour de 10.5. Il n'échange que lorsque cette valeur est en dehors de votre cotation, et en moyenne il prend alors l'écart qui dépasse $h$. Pour $h$ compris entre 0.5 et 1.5, ces écarts valent 0, $1.5 - h$ et $2.5 - h$, chacun deux fois sur six, si bien que la perte moyenne pour lui vaut $(4 - 2h)/3$.

Le pile ou face vous rapporte $h$ en moyenne, donc le point mort s'écrit
\[
(1 - \alpha)\, h = \alpha \cdot \frac{4 - 2h}{3}.
\]
Avec $\alpha = \tfrac12$, on obtient $h = 0.8$, valeur qui se situe bien dans l'intervalle supposé. Les deux autres cas se résolvent de la même façon, chacun avec son propre intervalle~: $\alpha = \tfrac14$ donne un $h$ inférieur à 0.5, où la perte vaut $(4.5 - 3h)/3$, et $\alpha = \tfrac34$ utilise à nouveau l'intervalle de 0.5 à 1.5.

\iqwrong{Utiliser l'écart type de la somme, soit 2.96. Un seul dé vu leur apprend bien moins que cela, et leur valeur ne varie que de $\pm 2.5$, si bien que la largeur nécessaire est beaucoup plus petite.}
\end{solution}

\begin{solution}{iq:iv2:making-a-market:23}
Les cotations sont données ci-dessous. La largeur passe pour la première fois sous 5 après quatre ventes. Elle se resserre à mesure que vous apprenez, alors que l'on pourrait s'attendre à ce qu'elle s'élargisse avec les pertes successives. Une série de quatre ventes a une probabilité de $1/256$, soit environ 0.4~\%, si la valeur vaut bien 100.
\begin{center}\small
\begin{tabular}{@{}lccccc@{}}
Ventes jusqu'ici & 0 & 1 & 2 & 3 & 4 \\ \hline
Cours acheteur & 25.00 & 10.00 & 3.57 & 1.22 & 0.41 \\
Cours vendeur & 75.00 & 50.00 & 25.00 & 10.00 & 3.57 \\
Largeur & 50.00 & 40.00 & 21.43 & 8.78 & 3.16
\end{tabular}
\end{center}

Une vente a une probabilité $0.75$ si la valeur vaut 0 et $0.25$ si elle vaut 100, si bien que chaque vente divise par 3 les chances en faveur de 100. Après $n$ ventes, ces chances valent $3^{-n}$. Le cours vendeur (ask) est la valeur après un achat supplémentaire, aux chances $3^{1-n}$, et le cours acheteur (bid) est la valeur après une vente supplémentaire, aux chances $3^{-n-1}$. Ainsi, le cours vendeur après $n$ ventes égale le cours acheteur après $n - 2$.

La largeur traduit ce qu'une seule transaction peut encore vous apprendre. Au départ, la valeur est entièrement inconnue. Après quatre ventes, vous êtes presque certain, si bien qu'une transaction n'a plus grand-chose à révéler.

\iqwrong{Penser que la largeur augmente lorsque le flux est unilatéral. Votre cours acheteur baisse vite, ce qui est pénible, mais la largeur est une autre chose. Ce qui élargit un vrai marché, c'est une incertitude nouvelle, et elle n'est pas présente dans ce modèle.}
\iqfollow{Et si la valeur vaut 0~?}{Alors la même série a une probabilité $(3/4)^4 \approx 32~\%$, ce qui n'a rien d'extraordinaire. C'est ce rapport qui vous fait y croire.}
\iqwhy{Après une forte chute accompagnée de ventes continues, un candidat qui élargit automatiquement a confondu la douleur avec l'incertitude.}
\end{solution}

\begin{solution}{iq:iv2:making-a-market:24}
Sur une échelle de 0 à 1, le cours vendeur vaut $a = \dfrac{(1 + \alpha) - \sqrt{1 - \alpha^2}}{2\alpha}$, et le cours acheteur vaut $1 - a$. Pour $\alpha = 0.5$, la cotation est de 36.6 au cours acheteur et 63.4 au cours vendeur. Lorsque $\alpha$ tend vers 1, le cours vendeur monte jusqu'à 100 et le cours acheteur descend jusqu'à 0, si bien qu'aucune transaction n'a lieu~: le marché se désorganise\index{désorganisation du marché}.

Le cours vendeur égale l'espérance conditionnelle à un achat. Un achat vient d'un trader informé lorsque $V > a$, et du pile ou face à toute valeur de $V$. Donc
\[
a = \frac{\alpha \int_a^1 v\,dv + \frac{1 - \alpha}{2} \int_0^1 v\,dv}{\alpha (1 - a) + \frac{1 - \alpha}{2}}
= \frac{\frac{\alpha}{2}(1 - a^2) + \frac{1 - \alpha}{4}}{\alpha(1 - a) + \frac{1 - \alpha}{2}}.
\]
En chassant les fractions, on obtient $\alpha a^2 - (1 + \alpha)a + \frac{1 + \alpha}{2} = 0$, et la plus petite racine est celle qui se trouve dans $[0, 1]$.

Pour $\alpha$ petit, $a \approx \frac12 + \frac{\alpha}{4}$, soit un cours vendeur d'environ $50 + 25\alpha$. Par exemple, $\alpha = 0.1$ donne 52.5, et $\alpha = 0.99$ donne 93.4.

\iqfollow{Pourquoi le marché ne se désorganise-t-il qu'à $\alpha = 1$ ici~?}{Parce qu'un flux non informé, quelle que soit sa taille, rémunère une cotation suffisamment large. Avec un continuum de valeurs, le profit d'un trader informé diminue à mesure que le cours vendeur se rapproche de 100, tandis que le flux conserve sa commission. Avec des coûts réels, la désorganisation peut survenir plus tôt.}
\iqwhy{La liquidité disparaît lorsque l'on craint un trading informé, comme c'est le cas autour d'une grande annonce inconnue. C'est le modèle le plus simple qui le montre.}
\end{solution}

\begin{solution}{iq:iv2:making-a-market:25}
La demi-fourchette optimale est $h = 250/21 \approx 11.9$, et elle rapporte environ 0.95 par trader. Vous n'atteignez le point mort qu'à $h = 50/7 \approx 7.14$. Au-delà d'une part informée de 25~\%, aucune cotation n'est rentable.

Un trader informé prend $(V - 50 - h)^+$ à la hausse, et autant à la baisse, si bien que, avec $V$ uniforme, cela vaut en moyenne $(50 - h)^2/100$. Un trader non informé échange avec la probabilité $1 - h/20$ et paie $h$ en moyenne. Avec une part $\alpha$,
\[
\pi(h) = (1 - \alpha)\, h \left(1 - \frac{h}{20}\right) - \alpha\, \frac{(50 - h)^2}{100}.
\]
Annuler la dérivée donne $h = 1\big/\left(\frac{1 - \alpha}{10} + \frac{\alpha}{50}\right)$. Pour $\alpha = 0.2$, cela vaut $1/0.084 \approx 11.9$. Poser $\pi(h) = 0$ donne $h = 50/7$ pour la plus petite racine.

Les deux cotations racontent donc des histoires différentes. Un teneur de marché sans concurrence prend 11.9. La concurrence fait descendre la demi-fourchette à 7.14, où le profit est nul. Pour $\alpha = \tfrac14$, la demi-fourchette optimale vaut 12.5, et elle rapporte exactement zéro, si bien qu'au-delà de cette part il ne faut pas coter du tout.

\iqwrong{Supposer que plus large est toujours mieux. C'est mieux face aux traders informés, mais vous perdez des clients, et c'est cela qui fixe la limite.}
\iqfollow{Pourquoi le marché ne se désorganise-t-il pas progressivement~?}{Parce que la commission du flux non informé a un plafond, alors que le gain des traders informés croît. Lorsque le gain maximal que vous pouvez réaliser passe sous zéro, vous vous arrêtez, et vous vous arrêtez d'un coup.}
\end{solution}

\begin{solution}{iq:iv2:making-a-market:26}
À plat, vous affichez 99.49 au cours acheteur, 100.51 au cours vendeur. Long de 40, vous affichez 98.69 au cours acheteur, 100.00 au cours vendeur. Le cours vendeur atteint 100 à partir d'une position de 26.

Le côté du client est un pile ou face, donc on choisit chaque prix séparément. Si une transaction modifie votre pénalité de $c$, vous maximisez $(\delta - c)\,e^{-2\delta}$, dont la dérivée s'annule en $\delta = \tfrac12 + c$.

À plat, chaque transaction coûte 0.01 de pénalité, donc $\delta = 0.51$. Long de 40, un achat coûte $0.01(41^2 - 40^2) = 0.81$, si bien que la profondeur à l'achat vaut $1.31$. Une vente économise $0.79$, donc $c = -0.79$, et la formule donne $\delta = -0.29$. Mais un cours vendeur en dessous de 100 n'obtient pas plus d'exécutions qu'un cours vendeur à 100, et ne fait que baisser le prix. On s'arrête donc à $\delta = 0$, et l'on gagne 0.79. C'est le cas limite du problème.

Votre cours vendeur atteint 100 lorsque $0.01(2q - 1) \ge \tfrac12$, soit à $q = 26$. Votre score espéré sur la transaction vaut environ 0.41, contre environ 0.18 pour la cotation à plat à la même position.

\iqwrong{Afficher une profondeur négative~: coter sous la juste valeur revient à offrir de l'argent, dès lors que la probabilité d'exécution a atteint un.}
\iqfollow{Que représente le $\tfrac12$~?}{Sans coût de stock, on maximise $\delta e^{-k\delta}$, ce qui donne $1/k$. C'est la marge qu'un vendeur faisant face à cette courbe de demande appliquerait.}
\end{solution}

\begin{solution}{iq:iv2:making-a-market:27}
Sans signal, on cote 97 au cours acheteur et 103 au cours vendeur. Avec le signal indiquant 110, on cote environ 103.66 au cours acheteur et 107.63 au cours vendeur, soit une largeur de 3.97 contre 6. Un signal exact à 60~\% seulement laisse environ 98.94 à 104.72, soit une largeur de 5.78.

Un signal public modifie votre a priori, et la même étape de Bayes donne ensuite la cotation. Partez du signal, de sorte que la probabilité de 110 vaut 0.8, et traitez un achat comme une observation supplémentaire, avec des vraisemblances $0.65$ et $0.35$. Alors
\[
\Pr(110 \mid \text{achat}) = \frac{0.8 \cdot 0.65}{0.8 \cdot 0.65 + 0.2 \cdot 0.35} \approx 0.881,
\]
ce qui donne un cours vendeur de $90 + 20 \times 0.881 = 107.63$. Une vente inverse les vraisemblances, et cela donne 103.66.

Notez que la cotation n'est pas centrée sur la nouvelle juste valeur, qui vaut $0.8 \times 110 + 0.2 \times 90 = 106$. Le mid se situe à 105.64, car une transaction peut vous rapprocher davantage de 90 que de 110.

\iqwrong{Penser que la part des traders informés a baissé. Elle n'a pas bougé. Mais la connaissance publique réduit ce qu'ils savent de plus que vous, si bien que chacune de leurs transactions apporte moins.}
\iqfollow{Que ferait un signal parfait~?}{Il ramènerait la largeur à zéro, puisque personne ne sait rien que vous ne sachiez. La valeur d'un signal pour un teneur de marché tient à une sélection adverse plus faible, ce qui signifie des cotations plus serrées et davantage de flux.}
\end{solution}

\begin{solution}{iq:iv2:making-a-market:28}
Pour $m = 0.5$, vous perdez environ 0.236 par transaction. Le seuil de rentabilité est $m = 35/36 \approx 0.97$. Avec trois teneurs de marché, il vaut $35/24 \approx 1.46$.

Vous ne vendez que lorsque votre propre dé est le plus bas. Votre cours vendeur vaut $d_1 + 3.5 + m$, et le contrat rapporte $d_1 + d_2$, donc votre gain sur une vente vaut $3.5 + m - d_2$, et gagner sélectionne un $d_2$ élevé. Si votre dé est le plus bas, vous gagnez tout, et en cas d'égalité vous gagnez la moitié. Ainsi, pour $d_2$ donné, votre part vaut $(d_2 - \tfrac12)/6$, et
\[
\text{gain} = \frac{1}{36} \sum_{d_2 = 1}^6 \left(3.5 + m - d_2\right) \left(d_2 - \tfrac12\right) = \frac{18m - 17.5}{36}.
\]
Ceci s'annule pour $m = 35/36$. Parmi les transactions où votre dé est strictement plus bas, le dé de l'adversaire vaut en moyenne $14/3$, et non 3.5. C'est la malédiction du vainqueur\index{malédiction du vainqueur}~: vous remportez la transaction lorsque vous avez sous-estimé.

Avec trois teneurs de marché, le même décompte sur 216 issues donne $m = 35/24$. Davantage d'adversaires aggravent la situation, car gagner en dit davantage.

\iqwrong{Utiliser la valeur inconditionnelle de la somme et ajouter un petit avantage. Cet avantage est entièrement absorbé par le fait que vous ne vendez que lorsque l'autre dé est élevé.}
\iqfollow{Est-ce la même chose que la sélection adverse liée à un client~?}{La structure est la même, mais le client ici ne sait rien. L'information vient de votre concurrent, à travers le choix du client entre vous deux.}
\end{solution}

\begin{solution}{iq:iv2:making-a-market:29}
À la première manche, coter 25 au cours acheteur et 75 au cours vendeur. Après un achat, coter 50 au cours acheteur et 83.33 au cours vendeur. Après un achat puis une vente, coter 50 au cours acheteur et 50 au cours vendeur. Si chaque manche fait intervenir une personne différente, ces trois mêmes cotations deviennent 25 à 75, 50 à 90, et 25 à 75.

Il y a quatre états~: la personne est informée ou non, et la valeur vaut 0 ou 100, chacun avec une probabilité de $\tfrac14$. Un achat a une vraisemblance de 1 venant d'une personne informée dont la valeur vaut 100, de 0 venant d'une personne informée dont la valeur vaut 0, et de $\tfrac12$ dans les autres cas. Ainsi, après un achat, les états ont pour probabilités~: personne informée avec une valeur de 100, $\tfrac12$~; personne non informée avec une valeur de 100, $\tfrac14$~; personne non informée avec une valeur de 0, $\tfrac14$. La valeur vaut 75.

Le cours vendeur suivant est la valeur conditionnée à un second achat, qui repondère ces probabilités en $\tfrac12 : \tfrac18 : \tfrac18$, d'où $100 \times \tfrac{5/8}{6/8} = 83.33$. Le cours acheteur suivant est la valeur conditionnée à une vente. Une personne informée qui vient d'acheter ne peut pas vendre, si bien qu'une vente prouve qu'elle ne sait rien, et la valeur revient exactement à 50. Comme une personne qui renverse sa position n'est pas informée et que ses transactions ne disent plus rien, aucune largeur n'est nécessaire.

\iqwrong{Traiter une même personne comme un nouveau tirage à chaque manche. Une deuxième transaction de la même personne est beaucoup moins informative qu'une transaction d'une personne nouvelle, si bien que le cours vendeur est inférieur à 90, et qu'un renversement est décisif.}
\iqfollow{Que se passe-t-il pour le cours acheteur si elle continue d'acheter~?}{Il reste à 50. Le cours vendeur passe de 83.33 à 90, puis vers 100, car chaque achat rend plus probable qu'elle sache.}
\end{solution}

\begin{solution}{iq:iv2:making-a-market:30}
La valeur inconditionnelle vaut $33/4 = 8.25$. La meilleure cotation est 7.5 au cours acheteur et 9.5 au cours vendeur, qui perd 0.05 en moyenne. Elle est centrée sur 8.5, au-dessus de la valeur. Si la personne achète, le contrat vaut 10, et si elle passe, il vaut $145/18 \approx 8.06$. Elle ne vend jamais.

Sa valeur ne dépend que de la carte qu'elle a vue~:
\begin{center}\small
\begin{tabular}{@{}l*{10}{c}@{}}
Carte & 1 & 2 & 3 & 4 & 5 & 6 & 7 & 8 & 9 & 10 \\ \hline
Valeur & 7.67 & 7.67 & 7.67 & 7.69 & 7.78 & 7.94 & 8.22 & 8.64 & 9.22 & 10.00
\end{tabular}
\end{center}
Les cartes basses comptent à peine, car l'une des deux autres est en général plus élevée. Les valeurs se resserrent en bas et s'écartent en haut.

Une cotation centrée sur 8.25, soit de 7.25 à 9.25, n'a pas non plus de vendeur, mais elle vend à la carte 10 au prix de 9.25, en perdant $0.75/10 = 0.075$. Remonter à 7.5 à 9.5 réduit cette perte à 0.05. À 7.75 à 9.75, les cartes 1 à 4 commencent à vous vendre, et la perte remonte à environ 0.056. Avec une distribution asymétrique de leur valeur, la meilleure cotation n'est pas centrée sur la moyenne. Pour vérifier, $0.9 \times 145/18 + 0.1 \times 10 = 8.25$.

\iqwrong{Centrer la cotation sur la valeur par habitude. Une cotation vous protège contre la valeur de l'autre partie, elle doit donc se situer autour de la distribution de cette valeur, qui n'est pas symétrique ici.}
\end{solution}

\begin{solution}{iq:iv2:making-a-market:31}
Si la personne a acheté $S$, alors $S$ vaut 9, et $D$ reste à 7. Si elle achète les deux, elles valent toutes deux 9, et elle a vu $P$. Si elle achète $S$ et vend $D$, alors $S$ vaut 9 et $D$ vaut 5, et elle a vu $Q$. Vous perdez 2 en moyenne par transaction.

Ayant vu $P = p$, elle valorise les deux contrats à $p + 3.5$. Ayant vu $Q = q$, elle valorise $S$ à $q + 3.5$ et $D$ à $10.5 - q$. Dans les deux cas, elle achète $S$ lorsque le dé vu vaut 5 ou 6. Ainsi, un achat de $S$ fait monter sa valeur de $5.5 - 3.5 = 2$, quel que soit le dé en cause. Pour $D$, il monte de 2 ou baisse de 2 avec des chances égales, donc il reste inchangé en moyenne.

Sa transaction sur $D$ vous indique alors quel dé elle a vu. Acheter les deux signifie que $P$ était le dé vu, donc $D$ monte à 9. Acheter $S$ et vendre $D$ signifie que c'était $Q$, donc $D$ baisse à 5.

Si le dé vu vaut 3 ou 4, elle passe sur les deux, ce qui arrive avec une probabilité de $\tfrac13$. Sinon, vous perdez 1.5 sur chacune des deux transactions, et $\tfrac23 \times 3 = 2$.

\iqwrong{Faire bouger $D$ avec $S$ parce qu'ils sont liés. Leur covariance est nulle, car ils ont la même variance issue de $P$ et de $Q$, avec des signes opposés pour $Q$, si bien qu'une transaction sur $S$ seule ne vous apprend rien sur $D$ en moyenne.}
\iqfollow{Alors comment actualiser $D$ après n'avoir vu que la transaction sur $S$~?}{Gardez le mid, mais élargissez la cotation~: $D$ vaudra désormais 5 ou 9, et vous ne savez pas encore lequel. Votre incertitude à son sujet a grandi alors que votre meilleure estimation n'a pas bougé.}
\end{solution}

\begin{solution}{iq:iv2:making-a-market:32}
Coter chaque fois au milieu, ce qui est une recherche dichotomique\index{recherche dichotomique}~: commencer à 5 au cours acheteur et 6 au cours vendeur, puis diviser ce qui reste. La perte totale espérée vaut $24/10 = 2.4$. Avec les nombres de 1 à 16, elle vaut $4.75$. Remonter depuis 1 au cours acheteur et 2 au cours vendeur coûte 7.

Commencez à 5 et 6. S'ils passent, le nombre vaut 5 ou 6, et vous ne perdez plus jamais, si bien que vous affichez 5 au cours acheteur et 6 au cours vendeur pour toujours. S'ils vendent, le nombre est compris entre 1 et 4, et vous avez perdu $5 - x$. Cotez ensuite 2 au cours acheteur et 3 au cours vendeur. Si le nombre vaut 1, ils vendent à nouveau et vous perdez 1, puis vous le savez. S'il vaut 4, ils achètent, et vous perdez 1. Les pertes totales pour les nombres de 1 à 10 sont donc~:
\begin{center}\small
\begin{tabular}{@{}l*{10}{c}@{}}
Nombre & 1 & 2 & 3 & 4 & 5 & 6 & 7 & 8 & 9 & 10 \\ \hline
Perte totale & 5 & 3 & 2 & 2 & 0 & 0 & 2 & 2 & 3 & 5
\end{tabular}
\end{center}
Cela fait 24 au total. La première cotation coûte 2.0 en moyenne, et la seconde seulement 0.4 de plus. Une programmation dynamique\index{programmation dynamique} portant sur tous les ensembles de nombres restants confirme qu'aucune autre cotation ne fait mieux, et que davantage de manches ne coûtent rien une fois le nombre connu.

Remonter vers le haut perd beaucoup plus. Si le nombre vaut 10, vous perdez 8, puis 6, 4 et 2, soit un total de 20, et la moyenne sur tous les nombres vaut 7. Chaque cotation se situe du mauvais côté du même nombre.

\iqfollow{Combien paieriez-vous pour qu'on vous révèle le nombre~?}{Jusqu'à 2.4, car c'est ce que coûte à vos cotations l'information qu'il faut découvrir.}
\iqwhy{Après une transaction contre vous, un déplacement agressif de la cotation coûte moins au total qu'une remontée progressive. Un candidat qui bouge un peu à chaque fois paie le plus cher.}
\end{solution}

\begin{solution}{iq:iv2:making-a-market:33}
Retirez un côté lorsqu'un lot supplémentaire ajouterait une pénalité plus lourde que ce qu'il peut rapporter. La position la plus grande que vous accepterez, selon le nombre de manches restantes de 1 à 10, est 2, 3, 3, 4, 4, 5, 5, 5, 6, 6. Le meilleur score espéré vaut 8.41, contre 8.00 sans jamais retirer de côté, et 8.35 pour la meilleure limite fixe, qui est $\pm 3$.

Remontez depuis la fin par induction à rebours\index{induction à rebours}. Avec une seule manche restante, un achat à la position $q$ rapporte $1 - 0.2(2q + 1)$, ce qui est positif pour $q = 1$ et exactement nul pour $q = 2$. Vous achetez donc jusqu'à 2 et pas davantage. Avec deux manches restantes, calculez la valeur de chaque position avec une manche restante, puis comparez à nouveau. Un achat maintenant vaut plus qu'il ne vaudrait à la fin, car une transaction supplémentaire peut vous ramener en arrière. Vous achetez donc jusqu'à 3.

Chaque pas en arrière ajoute du temps pour que le flux vous ramène vers une position neutre, si bien que la position autorisée croît avec le nombre de manches restantes.

Ne jamais retirer de côté rapporte un par manche, soit 10, moins $0.2 \times \mathbb{E}[q^2] = 0.2 \times 10$, soit 8.00. Une limite fixe de $\pm 3$ capte l'essentiel du gain, mais elle est trop serrée au début et trop large à la fin.

\iqwrong{Traiter une limite comme un nombre fixe. Elle doit dépendre du temps dont vous disposez pour sortir de la position. Une position qui a une journée pour être débouclée n'est pas la même qu'une position à la clôture.}
\iqfollow{Et si les clients réagissent au prix~?}{Vous pouvez alors incliner votre cotation au lieu de refuser des transactions. Retirer un côté est une inclinaison extrême, et un prix fait le même travail de façon plus souple.}
\end{solution}

\begin{solution}{iq:iv2:making-a-market:34}
$\lambda = \sigma_v/(2\sigma_u) = 0.01$ par action. L'ordre de l'initié vaut $x = (\sigma_u/\sigma_v)(V - 100) = 50(V - 100)$, et le gain espéré vaut $\sigma_v\sigma_u/2 = 100$. Une fois le cours fixé, la moitié de la variance de $V$ subsiste, soit 2 sur 4. C'est le modèle de Kyle à une période\index{modèle de Kyle}.

L'initié choisit $x$ pour maximiser $x(V - 100 - \lambda x)$, puisque le bruit est de moyenne nulle. On obtient $x = \beta(V - 100)$ avec $\beta = 1/(2\lambda)$. Vous fixez le cours à l'espérance de la valeur conditionnellement au flux, si bien que $\lambda$ est la pente de la régression de $V$ sur $y$~:
\[
\lambda = \frac{\mathrm{Cov}(V, y)}{\mathrm{Var}(y)} = \frac{\beta \sigma_v^2}{\beta^2 \sigma_v^2 + \sigma_u^2}.
\]
En remplaçant $\beta = 1/(2\lambda)$, on obtient $4\lambda^2\sigma_u^2 = \sigma_v^2$, qui est le résultat.

L'ordre de l'initié a la même variance que le bruit, car $\beta\sigma_v = \sigma_u$. Il se fond exactement dans la foule. La moitié de la variance de $y$ est donc du signal, et le cours se déplace à mi-chemin de la vérité. Le gain espéré de l'initié vaut $\mathbb{E}\left[\beta (V-100)^2\right]/2 = 50 \times 4/2 = 100$, et ce sont les traders de bruit qui le paient~: chaque ordre $u$ fait bouger le cours contre lui-même de $\lambda u$, ce qui leur coûte $\lambda\sigma_u^2 = 100$ en moyenne.

\iqwrong{Faire trader l'initié jusqu'à ce que le cours atteigne la vraie valeur. L'ordre fait bouger le cours, si bien que la meilleure taille n'est que la moitié de celle qui ferait disparaître l'avantage.}
\iqfollow{Que se passe-t-il si le bruit double~?}{$\lambda$ est divisé par deux, l'ordre et le gain de l'initié doublent, et le cours n'est pas moins informatif. Davantage de bruit rémunère l'initié sans modifier la variance restante.}
\end{solution}

\begin{solution}{iq:iv2:making-a-market:35}
Le cours acheteur vaut environ 98.35, et le cours vendeur environ 100.05. La largeur vaut $\gamma\sigma^2 + \frac{2}{\gamma}\ln\left(1 + \frac{\gamma}{k}\right) \approx 0.40 + 1.29 = 1.69$, et le centre vaut $100 - \gamma\sigma^2 q = 99.2$, avec $\gamma = 0.1$, $\sigma^2 = 4$, $k = 1.5$ et $q = 2$. Lorsque $\gamma$ tend vers 0, chaque côté tend vers $1/k$ à partir de 100, soit environ 99.33 à 100.67.

\begin{sloppypar}Sans exécution, l'utilité exponentielle d'une loi normale donne $U_0 = -\exp\left(-\gamma(x + 100q) + \tfrac12\gamma^2\sigma^2 q^2\right)$, où $x$ est votre liquidité. Si votre cours acheteur à la profondeur $\delta$ est exécuté, vous payez $100 - \delta$ et détenez $q + 1$, donc\end{sloppypar}
\[
U = U_0 \exp\left(-\gamma(\delta - d_b)\right), \qquad d_b = \tfrac12\gamma\sigma^2(2q + 1).
\]
Ici $d_b$ est le coût d'inventaire de cet achat. L'utilité espérée de ce côté vaut
\[
U_0\left[1 + e^{-k\delta}\left(e^{-\gamma(\delta - d_b)} - 1\right)\right].
\]
Comme $U_0$ est négative, vous voulez que le crochet soit petit. En posant $u = \delta - d_b$, la condition du premier ordre donne $e^{-\gamma u} = k/(k + \gamma)$, soit
\[
\delta = d_b + \frac{1}{\gamma}\ln\left(1 + \frac{\gamma}{k}\right).
\]
Pour le cours vendeur, les mêmes étapes donnent $d_a = \tfrac12\gamma\sigma^2(1 - 2q)$. Ainsi $\delta_b = 1.0 + 0.645$ et $\delta_a = -0.6 + 0.645 = 0.045$. En détenant 2 actions, vous facturez à peine la vente et vous facturez beaucoup l'achat.

C'est la forme à une période de la cotation d'Avellaneda--Stoikov\index{modèle d'Avellaneda--Stoikov}. Le décalage dépend de votre position et du risque, et la largeur dépend du risque et de la sensibilité des clients au prix.

\iqwrong{Ajouter le terme d'inventaire d'un seul côté. Il déplace les deux côtés, et $d_b + d_a = \gamma\sigma^2$ est la part qu'il ajoute à la largeur.}
\end{solution}
